The sum of the reciprocals of the squares is $\pi^2/6$. The Gaussian integral is $\sqrt\pi$. In the sense of the functional equation, $1+2+3+\cdots$ equals $-\frac1{12}$. Stirling's series and Gauss's arithmetic--geometric mean are approximations with a remainder one can name; the mean doubles the number of correct digits of $\pi$ at each step. This book computes those values. Its tools are finite calculus and generating functions, the Gamma function, Fourier analysis on the circle, on $\Z/n\Z$ (where it becomes the fast Fourier transform) and on the line, residues, and the zeta function with its functional equation. One chapter spends them on a single question, the number of ways to write $n$ as a sum of positive integers, and proves the formula of Hardy and Ramanujan. Every theorem is proved. A few results are quoted and marked, among them Rademacher's series for the number of partitions; nothing later depends on them. The exercises are graded, and each has a complete solution.
